% !PRINTED 301
\weberpage{301}

\[
2\varrho^5 - \frac{A}{3}\varrho^4 - \frac{5B}{3}\varrho^3 + \frac{AB+5C}{3}\varrho^2 - \frac{4AC}{9}\varrho - \left(\frac{AB^2}{12} + \frac{13}{18}BC - \frac{D}{2}\right) = 0.
\]

Die Discriminante $P$ der quadratischen Form $p$ ist also von einer Gleichung fünften Grades abhängig, deren Coefficienten die Invarianten der ersten Elementarcovariante sind. In ähnlicher Weise würde man finden, dass für zwei Formen dritten Grades $f$ und $\varphi$ die zweite Elementarcovariante, welche in diesem Falle bereits Invariante ist, durch eine quadratische Gleichung sich aus den Coefficienten der ersten Elementarcovariante berechnen lässt.

277. \textit{Schlussbemerkungen.} Ich will die letzten Betrachtungen mit dem Hinweis auf deren Verallgemeinerung schliessen. Zunächst hatte Brill (Math. Annalen Bd. XX) die Bedeutung dieser ersten Elementarcombinante für alle übrigen Combinanten bemerkt, auch für den Fall, dass das zu Grunde gelegte simultane System aus mehr als zwei Formen $f, \varphi$ besteht. Unter anderm zeigte er, dass sie eine allgemeine Form ihres Grades ist und dass ihre Discriminante in zwei rationale Factoren zerfällt. Stephanos, der in den Comptes Rendus 1882 gleichzeitig auf diese Beziehungen aufmerksam machte, hat alsdann in seinen Abhandlungen: \glqq Mémoire sur les faisceaux de formes binaires ayant une même Jacobienne\grqq{} vom Jahre 1883, und \glqq Mémoire sur la Théorie des Formes binaires et sur l'élimination\grqq{} vom Jahre 1884 dieselben noch eingehender untersucht. Er sowohl als Franz Meyer (Apolarität und rationale Curven, Seite 391) zeigen, dass die in unserm speciellen Falle aufgestellte Endgleichung fünften Grades für die Coefficienten der zweiten Elementarcovariante im allgemeinen Falle, wenn $f$ und $\varphi$ vom $n^{\text{ten}}$ Grade sind, den Grad $N = \dfrac{(2(n-1))!}{(n-1)!\,n!}$ besitzt. Bezüglich weiterer interessanter Resultate müssen wir auf die citirten Quellen verweisen. Hier sei nur noch auf die Umkehrung der im Vorigen behandelten Aufgabe aufmerksam gemacht. Man kann sich nämlich fragen: welche Formen $f$ und $\varphi$ gehören zu einer gegebenen ersten Elementarcovariante $F$ von $f$ und $\varphi$? Für unser Beispiel lässt sich die Frage einfach beantworten. Es war $f = r_x^4$, $\varphi = s_x^4$, also:

\[
\psi = r_x^4 s_y^4 - s_x^4 r_y^4 = \mathrm{Const.}\, F_{y^3}(yx) + \mathrm{Const.}\, \Phi_y(yx).
\]

Da $F$ gegeben ist, $\Phi$ durch eine Gleichung fünften Grades aus $F$ sich berechnen lässt, $y$ aber willkürlich ist, so erkennt man, dass zu einer gegebenen Form $F$ stets fünf Schaaren oder lineare Systeme

% checked against the GDZ scan (27 Oct 2026): the capital Psi of this fount (arms curled outward), \Psi
\[
\Psi = \mu\, r_x^4 + \lambda\, s_x^4
\]

gehören. Einer solchen Schaar gehören insbesondere zwei specielle Formen an, für welche die Invariante $i = 0$ ist, also die beiden
%CONT>>


% ===== pdf idx 316, printed 302.0 =====
% !PRINTED 302
\weberpage{302}
%<<CONT
Tetraeder der Schaar, welche durch die aus $i = 0$ sich ergebende quadratische Gleichung für $\dfrac{\lambda}{\mu}$ bestimmt sind. In einer Mittheilung an die Société mathématique de France (1880) hat Stephanos gezeigt, dass, wenn $f$ und $\varphi$ Formen $n^{\text{ten}}$ Grades sind, die Zahl der linearen Systeme, welche zu einer gegebenen Functionaldeterminante zweier Formen $f$ und $\varphi$ gehören, gleich der oben angeführten Zahl $N = \dfrac{(2(n-1))!}{(n-1)!\,n!}$ ist.

\begin{center}
\S~29.\\
\textbf{Typische Darstellung der Form sechsten Grades.}
\end{center}

278. \textit{Allgemeiner Fall.} In den Systemen von Formen gerader Ordnung treten im allgemeinen Falle lineare Covarianten nicht auf. Dagegen besitzen dieselben im Allgemeinen, wenn $n > 4$, stets mindestens zwei quadratische Covarianten $\varphi$ und $\psi$, deren Resultante $R_{\varphi,\psi}$ nicht verschwindet. (Vergl. Clebsch, \glqq Binäre Formen\grqq, § 102, Seite 140.) Diese zwei quadratischen Formen können nun unter anderm als Grundlage für die typische Darstellung von Formen gerader Ordnung benutzt werden, wie Clebsch und Gordan, Annali di Mat., Ser. II., t. 1, gezeigt haben. Denn zwischen vier quadratischen Formen $a_x^2, l_x^2, m_x^2, n_x^2$ besteht stets die Relation (vergl. Nr. 135, I, worin $y = a$ zu setzen ist)

\[
a_x^2 \cdot R_{lmn} = l(a\lambda)^2 + m(a\mu)^2 + n(a\nu)^2,
\]

wobei $\lambda, \mu, \nu$ die Functionaldeterminanten $(mn)\,m_xn_x$, respective $(nl)\,n_xl_x$, $(lm)\,l_xm_x$ repräsentiren. Und diese Relation findet hier dieselbe Verwendung, wie bei den Formen ungeraden Grades die Identität

\[
a_x(\alpha\beta) = \alpha_x(a\beta) - \beta_x(a\alpha).
\]

Wir erheben sie auf die $\dfrac{n}{2}^{\text{te}}$ Potenz und erhalten die Beziehung

\weq{I}{R_{lmn}^{\frac{n}{2}} \cdot f = \left\{ l(a\lambda)^2 + m(a\mu)^2 + n(a\nu)^2 \right\}^{\frac{n}{2}},}

aus welcher die typische Darstellung von $f$ unmittelbar hervorgeht. Es treten dabei allerdings an Stelle der zwei alten Variabeln $x_1x_2$ drei neue $l, m, n$ auf; doch sind diese nicht von einander unabhängig. Denn zwischen ihnen besteht die Relation [vergl. Nr. 134, 3, $(\gamma)$]

\[
0 = \begin{vmatrix}
A_{ll}, & A_{lm}, & A_{ln}, & l \\
A_{lm}, & A_{mm}, & A_{mn}, & m \\
A_{ln}, & A_{mn}, & A_{nn}, & n \\
l, & m, & n, & 0
\end{vmatrix}
\]

und somit ist die typische Darstellung (I) von $f$ in der That eine binäre.


% ===== pdf idx 317, printed 303.0 =====
% !PRINTED 303
\weberpage{303}

279. Was nun die Form $f = a_x^6$ betrifft, so könnten wir deren typische Darstellung unmittelbar auf Gleichung (I) basiren. Indessen bietet in diesem Falle die in Nr. 135, (VII) gegebene Relation ein noch bequemeres Mittel, zum nämlichen Resultate zu gelangen, so lange die Invariante $R$ von $0$ verschieden ist.

Ersetzen wir in derselben nämlich $f, \varphi, \psi$ durch die drei quadratischen Formen $l, m, n$ des Systemes von $f = a_x^6$, so geht dieselbe über in:

% checked against the GDZ scan (27 Oct 2026): the mark above the subscript of A_{mm} is an ink speck; plain A_{mm}
\[
\begin{vmatrix}
A_{ll}, & A_{lm}, & A_{ln}, & l_x^2, & l_y^2 \\
A_{lm}, & A_{mm}, & A_{mn}, & m_x^2, & m_y^2 \\
A_{ln}, & A_{mn}, & A_{nn}, & n_x^2, & n_y^2 \\
l_x^2, & m_x^2, & n_x^2, & 0, & (xy)^2 \\
l_y^2, & m_y^2, & n_y^2, & (xy)^2, & 0
\end{vmatrix} = 0.
\]

Entwickeln wir diese Determinante mit Hilfe des Additionstheorems für Determinanten (vergl. Bd. I Nr. 35 u. 36) nach Potenzen von $(xy)$, so erhalten wir:

\[
0 = -2R^2\cdot(yx)^4 - 2
\begin{vmatrix}
A_{ll}, & A_{lm}, & A_{ln}, & l_x^2 \\
A_{lm}, & A_{mm}, & A_{mn}, & m_x^2 \\
A_{ln}, & A_{mn}, & A_{nn}, & n_x^2 \\
l_y^2, & m_y^2, & n_y^2, & 0
\end{vmatrix}(yx)^2 +
\begin{vmatrix}
A_{ll}, & A_{lm}, & A_{ln}, & l_x^2, & l_y^2 \\
A_{lm}, & A_{mm}, & A_{mn}, & m_x^2, & n_y^2 \\
A_{ln}, & A_{mn}, & A_{nm}, & n_x^2, & n_y^2 \\
l_x^2, & m_x^2, & n_x^2, & 0, & 0 \\
l_y^2, & m_y^2, & n_y^2, & 0, & 0
\end{vmatrix}.
\]
% [sic] third row, third entry printed "A_{nm}" where the pattern elsewhere in this row/column uses "A_{mn}"

Hierin besitzt der Coefficient von $2(yx)^2$ gemäss der Relation (VI) Nr. 135 den Werth $-2R^2(yx)^2$, wie sich aus derselben wiederum nach dem Additionsgesetze für Determinanten ergiebt. Demnach geht die letzte Gleichung, wenn wir überdies $y$ durch $a$ ersetzen und mit $a_x^2$ multipliciren, über in:

\[
-2R^2\cdot f = a_x^2
\begin{vmatrix}
A_{ll}, & A_{lm}, & A_{ln}, & l, & (al)^2 \\
A_{lm}, & A_{mm}, & A_{mn}, & m, & (am)^2 \\
A_{ln}, & A_{mn}, & A_{nn}, & n, & (an)^2 \\
l, & m, & n, & 0, & 0 \\
(al)^2, & (am)^2, & (an)^2, & 0, & 0
\end{vmatrix}.
\]

280. Entwickelt man diese Determinante nach den Minoren der letzten zwei Zeilen und Colonnen (vergl. insbesondere Bd. I Seite 72, Beispiel 2), so erhält man:

\weq{I}{\begin{aligned}
-2R^2\cdot f ={}& \sum A_{nn}\left\{ l^2(f,m^2)^4 - 2lm(f,ml)^4 + m^2(f,l^2)^4 \right\} \\
&{}- 2\sum A_{mn}\left\{ l^2(f,mn)^4 - lm(f,ln)^4 + mn(f,l^2)^4 - nl(f,ml)^4 \right\},
\end{aligned}}
% [sic] this display is printed with the margin label "(I)", reusing the same numeral already used for the equation on p. 302
%CONT>>


% ===== pdf idx 318, printed 304.0 =====
% !PRINTED 304
\weberpage{304}
%<<CONT
wobei jede Summe sich über drei Glieder erstreckt, deren zweites und drittes aus den angeschriebenen durch cyclische Vertauschung von $l\, m\, n$ hervorgehen.

Die in der Entwicklung rechts auftretenden quadratischen Formen

\[
(f, l^2)^4,\quad (f, m^2)^4,\quad (f, n^2)^4,\quad (f, lm)^4,\quad (f, ln)^4,\quad (f, mn)^4
\]

sind aber, wie wir bereits gezeigt haben, lineare Functionen der Covarianten $l, m, n$. Diese Entwicklung rechts ist also eine Function dritten Grades in den Variabeln $l, m, n$, und die darin auftretenden Coefficienten sind Invarianten. Die typische Darstellung von $f$ ist also in der That durch die Gleichung (I) geleistet. Betrachtet man in ihr die Covarianten $l, m, n$ als drei Coordinaten $x_1, x_2, x_3$, so stellt die rechte Seite dieser Gleichung eine Curve dritter Ordnung dar. Da aber nach der Schlussbemerkung in Nr. 278 dieselben drei Coordinaten den Kegelschnitt

% checked against the GDZ scan (27 Oct 2026): the third row prints lowercase l_{ln}, l_{mn} (misprint for A_{ln}, A_{mn}); reproduced as printed
\weq{II}{
\begin{vmatrix}
A_{ll} & A_{lm} & A_{ln} & x_1 \\
A_{lm} & A_{mm} & A_{mn} & x_2 \\
l_{ln} & l_{mn} & A_{nn} & x_3 \\
x_1 & x_2 & x_3 & 0
\end{vmatrix} = 0
}

befriedigen, so erkennt man: Die Wurzeln der Gleichung sechsten Grades $f = 0$ werden durch die typische Darstellung von $f$ als Schnittpunkte einer Curve dritter Ordnung und eines Kegelschnittes erhalten.

281. \textit{Der specielle Fall $R = 0$.} Wir hatten bisher vorausgesetzt, dass $R$ von null verschieden sei. Nehmen wir nun an, dass $R$ identisch verschwinde, so dürfen, wenn die typische Darstellung gleichwohl möglich sein soll, die Covarianten $l, m, n$ keinen gemeinsamen Factor haben. Zu dem Zwecke genügt es, dass $A_{ll}A_{mm} - A_{ml}^2$ von null verschieden sei.

Wir benutzen alsdann behufs typischer Darstellung die in der Theorie der quadratischen Formen gegebene Identität [vgl. Nr. 135, (V)]

\[
2v_x^2 \cdot v_y^2 =
\begin{vmatrix}
A_{ll} & A_{lm} & l \\
A_{lm} & A_{mm} & m \\
l_y^2 & m_y^2 & (xy)^2
\end{vmatrix},
\]

welche, nach $(xy)^2$ entwickelt, die Form annimmt:

\[
(xy)^2 A_{vv} = 2v_x^2v_y^2 -
\begin{vmatrix}
A_{ll} & l_{lm} & l \\
A_{lm} & A_{mm} & m \\
l_y^2 & m_y^2 & 0
\end{vmatrix}.
\]
% [sic] second row/first row printed "l_{lm}" where the pattern elsewhere suggests "A_{lm}"

Erheben wir dieselbe auf die dritte Potenz und ersetzen alsdann $y$ durch $a$, so kommt:
%CONT>>


% ===== pdf idx 319, printed 305.0 =====
% !PRINTED 305
\weberpage{305}
\[
\begin{aligned}
A_{\nu\nu}^3\cdot f&=\prod_{i=1}^{i=3}\left\{2\,\nu\,(\nu_ia)^2-\begin{vmatrix}A_{ll}, & A_{lm}, & l\\ A_{lm}, & A_{mm}, & m\\ (l_ia)^2, & (m_ia)^2, & 0\end{vmatrix}\right\}\\
&=8\,\nu^3\cdot(f,\,\nu^3)^6-4\,\nu^2\sum{}'(\nu_1a)^2(\nu_2a)^2\begin{vmatrix}A_{ll}, & A_{lm}, & l\\ A_{lm}, & A_{mm}, & m\\ (l_3a)^2, & (m_3a)^2, & 0\end{vmatrix}\\
&\quad+2\,\nu\sum{}'(\nu_1a)^2\begin{vmatrix}A_{ll}, & A_{lm}, & l\\ A_{lm}, & A_{mm}, & m\\ (l_2a)^2, & (m_2a)^2, & 0\end{vmatrix}\cdot\begin{vmatrix}A_{ll}, & A_{lm}, & l\\ A_{lm}, & A_{mm}, & m\\ (l_3a)^2, & (m_3a)^2, & 0\end{vmatrix}\\
&\quad-\prod_{i=1}^{i=3}\begin{vmatrix}A_{ll}, & A_{lm}, & l\\ A_{lm}, & A_{mm}, & m\\ (l_ia)^2, & (l_ia)^2, & 0\end{vmatrix},% [sic] printed (l_ia)^2 twice in this determinant, not (l_ia)^2,(m_ia)^2
\end{aligned}
\]
wobei sich die beiden Summen $\sum{}'$ über je drei Glieder erstrecken.

In der Entwicklung rechts verschwindet das erste und dritte Glied. Denn die Invarianten
\[
(f,\,\nu^3)^6,\quad (f,\,\nu l^2)^6,\quad (f,\,\nu m^2)^6,\quad (f,\,\nu lm)^6,
\]
welche in diesen zwei Gliedern als Coefficienten auftreten, sind schiefe Invarianten von den Graden 25, bezw. 15, 19, 17 in den Coefficienten von $f$, sie müssen daher, da nur die einzige schiefe Fundamentalinvariante $R$ im Systeme existirt, diese zum Factor haben, und sonach gleichzeitig mit $R$ verschwinden.

Es reducirt sich also die typische Darstellung auf:
\weq{III}{A_{\nu\nu}^3 f\equiv -\prod\begin{vmatrix}A_{ll}, & A_{lm}, & l\\ A_{lm}, & A_{mm}, & m\\ (l_ia)^2, & (l_ia)^2, & 0\end{vmatrix}% [sic] printed (l_ia)^2 twice, as above
-4\,\nu^2\sum{}'(\nu_1a)^2(\nu_2a)^2\begin{vmatrix}A_{ll}, & A_{lm}, & l\\ A_{lm}, & A_{mm}, & m\\ (l_3a)^2, & (m_3a)^2, & 0\end{vmatrix}.}

Da $\nu$ als Functionaldeterminante von $l$ und $m$ der Relation genügt
\weq{IV}{-2\,\nu^2=\begin{vmatrix}A_{ll}, & A_{lm}, & l\\ A_{lm}, & A_{mm}, & m\\ l, & m, & 0\end{vmatrix},}
so erhält man nach Substitution dieses Werthes von $\nu^2$ in die drei letzten Glieder der Gleichung (III) auf der rechten Seite einen in den Formen $l$ und $m$ homogenen cubischen Ausdruck von der Gestalt:
\[
A_{\nu\nu}^3 f=C_1l^3+C_2l^2m+C_3lm^2+C_4m^3.
\]


% ===== pdf idx 320, printed 306.0 =====
% !PRINTED 306
\weberpage{306}
Die Coefficienten $C_i$ dieser Darstellung sind Functionen der Invarianten $A_{ll}$, $A_{lm}$, $A_{mm}$ und ausserdem der Invarianten
\[
(f,\,l^3)^6,\quad (f,\,m^3)^6,\quad (f,\,l^2m)^6,\quad (f,\,lm^2)^6;
\]
denn die zwei Invarianten $(f,\,\nu^2l)^6$, $(f,\,\nu^2m)^6$ lassen sich sofort wegen (IV) auf dieselben reduciren.

Von diesen wurde die erste bereits früher berechnet (Nr. 264):
\[
(f,\,l^3)^6=2A_{ln}-\frac{2}{3}BA_{ll}+\frac{1}{3}A_{lm}.% [sic] index \textquotedbl{}ln\textquotedbl{} as printed
\]
Die drei andern ergeben sich direct aus den beiden ersten Relationen in (3) Nr. 266, indem man dieselben zweimal über $m$, resp. $l$ schiebt. Die Coefficienten der typischen Darstellung
\weq{V}{A_{\nu\nu}^3\cdot f=C_1l^3+C_2l^2m+C_3lm^2+C_4m^3}
sind somit vollständig bekannte Grössen.

Weitere specielle Fälle sind in Bezug auf typische Darstellung noch wenig untersucht worden. In den \glqq Binären Formen\grqq{} Clebsch's findet sich noch neben dem eben besprochenen Fall die typische Darstellung von $f$, wenn $A=0$, $B=0$, $C=0$, $D$ und $R$ von Null verschieden. Es ergiebt sich das Resultat:
\[
D^2f=2n^3+\frac{3D}{2}l^2m,
\]
wobei $n$ und $l$ reine Quadrate sind. (Vergl. a. a. O. Seite 467.)

\begin{center}
\S\ 30. \textbf{Auflösung specieller Gleichungen sechsten Grades.}
\end{center}

282. \textit{Ueberblick.} Von den Gleichungen sechsten Grades, welche vermöge gewisser Invarianten- oder Covariantenbedingungen den Charakter von speciellen Gleichungen tragen, sind einmal jene hervorzuheben, deren Invariantenrelationen ihre Reduction auf Gleichungen niedrigerern Grades involviren, dann aber auch jene, durch deren Invariantenbeziehungen zwar nicht ihr Grad reducirt wird, welche aber durch andere bemerkenswerthe Eigenschaften unser Interesse erregen, insofern sie mit Modulargleichungen oder Multiplicatorgleichungen identificirt werden können. Wir wollen einige wichtige Fälle dieser speciellen Gleichungen sechsten Grades im Folgenden eingehender untersuchen, insbesondere um auf die Art und Weise hinzudeuten, wie die hier einschlägigen Fragen mit unsern Mitteln behandelt werden können. Der uns zunächst interessirende Fall ist der Fall, in welchem $R=0$, und die Resultante $R_{lm}=-2A_{\nu\nu}$ von $l$ und $m$ von Null verschieden ist. Hieran schliessen sich die Untersuchungen über den Fall $R=0$, $A_{\nu\nu}=0$, die zu drei Unterabtheilungen von Bedingungen Veranlassung geben, jenachdem das Verschwinden von $A_{\nu\nu}$ eine Folge davon, dass $l$ und
%CONT>>


% ===== pdf idx 321, printed 307.0 =====
% !PRINTED 307
\weberpage{307}
%<<CONT
$m$ einen linearen Factor gemeinschaftlich besitzen, oder dass $l$ proportional $m$ ist, oder endlich, dass $l$ oder $m$ selbst identisch null wird. Im letzten Falle wird sich zeigen, dass, wenn eine der drei quadratischen Formen verschwindet, nothwendig auch die zwei andern null werden, und die Form $(f,f)^4$ entweder das Quadrat oder die vierte Potenz einer Form zweiten resp. ersten Grades ist, so dass also nur noch diese beiden Möglichkeiten zu besprechen sind, um den Fall $R=0$, $A_{\nu\nu}=0$ völlig zu erschöpfen.

Zum Schlusse weisen wir noch auf die beiden Fälle: $A=0$, $C=0$ und $B=\text{Const.}\times A^2$, $C=\text{Const.}\times A^3$ hin, unter welchen Bedingungen sich die allgemeine Gleichung sechsten Grades durch lineare Transformation auf Gleichungen mit einem einzigen Parameter reduciren lässt.

283. \textit{Erster Fall: $R=0$, $A_{\nu\nu}\gtrless 0$.} Die typische Darstellung der Form $f=a_x^6$ für $R=0$ am Schlusse des letzten Paragraphen lehrt:

\glqq Verschwindet die Invariante $R$, so ist die Gleichung sechsten Grades $f=0$ algebraisch lösbar, und ihre Wurzeln sind drei Elementenpaare einer Involution, deren Doppelpunkte durch die linearen Factoren von $\nu=0$ bestimmt sind.\grqq{}

Denn durch die Substitution
\[
y=\frac{l}{m}
\]
geht die unter Nr. 281 gegebene Gleichung (V) über in
\[
C_1y^3+C_2y^2+C_3y+C_4=0.
\]
Die Wurzeln dieser Gleichung seien $y_1$, $y_2$, $y_3$; die Wurzeln von $f=0$ werden alsdann durch die drei quadratischen Gleichungen:
\weq{1}{l-y_1m=0,\quad l-y_2m=0,\quad l-y_3m=0}
geliefert. Da nach Voraussetzung $A_{\nu\nu}\gtrless 0$, also $l$ und $m$ keinen gemeinschaftlichen Factor haben, so kann man diese beiden quadratischen Covarianten auf die canonische Form:
\[
\begin{aligned}
l&=\alpha_1\xi^2+\beta_1\eta^2\\
m&=\alpha_2\xi^2+\beta_2\eta^2
\end{aligned}
\]
durch eine gemeinschaftliche lineare Transformation bringen (vergl. Nr. 130), wobei $\xi$ und $\eta$ die linearen Factoren von $\nu=0$ sind. Alsdann gehen die drei quadratischen Formen (1) über in:
\[
c_1\xi^2-c_2\eta^2=0,\quad c_1'\xi^2-c_2'\eta^2=0,\quad c_1''\xi^2-c_2''\eta^2=0
\]
und diese drei Gleichungen repräsentiren drei Punktepaare, welche zu den Punkten $\xi=0$, $\eta=0$ harmonisch liegen, was zu beweisen war.


% ===== pdf idx 322, printed 308.0 =====
% !PRINTED 308
\weberpage{308}
284. \textit{Zweiter Fall; erste Unterabtheilung: $R=0$, $A_{\nu\nu}=0$.} Wir nehmen an, das Verschwinden von $A_{\nu\nu}=A_{ll}A_{mm}-A_{lm}^2$ sei dadurch bedingt, dass $l$ und $m$ einen gemeinschaftlichen linearen Factor $\alpha$ haben. Da aber die Bedingung $R=0$ eine lineare Relation:
\[
2\,\nu\cdot R=\begin{vmatrix}A_{ll}, & A_{lm}, & A_{ln}\\ A_{ml}, & A_{mm}, & A_{mn}\\ l, & m, & n\end{vmatrix}=c_1l+c_2m+c_2n=0% [sic] third coefficient printed c_2 twice, not c_3
\quad\text{(vgl. Nr. 135, IV$_{\text{b}}$)}
\]
zwischen den drei Covarianten $l$, $m$, $n$ involvirt, so muss auch $n$ mit $l$ und $m$ denselben linearen Factor $\alpha$ gemeinschaftlich haben.

Die Bedingungen für unsern Fall lassen sich demnach auch in der Form schreiben
\weq{1}{l=\alpha\cdot r,\quad m=\alpha\cdot s,\quad n=\alpha\cdot t,}
wobei die linearen Factoren $r$, $s$, $t$ von einander verschieden sein sollen. Aus diesen drei Bedingungen (1) gehen direct die Relationen
\weq{2}{(l\alpha)^2=0,\quad (m\alpha)^2=0,\quad (n\alpha)^2=0}
hervor, während nach der Theorie der quadratischen Formen
\[
\alpha^2=(l,\,m)=\nu\quad\text{ist.}
\]
Nun ist aber die Covariante $l$ nichts anderes als $(f,k)^4$ und es liegt die Frage nahe, ob eine dieser beiden Formen $f$ und $k$ den rationalen linearen Factor $\alpha_x$ besitzt. Um dies zu untersuchen, hat man nur die Ueberschiebungen von $k$, resp. $f$ über $\alpha$ zu bilden und nachzusehen, welche derselben identisch verschwindet. Man findet alsdann:
\weq{3}{(k\alpha)^3k_x=0}
und
\weq{4}{(f\alpha)^4f_x^2=0.}

Die erste Gleichung lehrt: $k$ besitzt den Factor $\alpha_x$ quadratisch, die zweite: $f$ enthält ihn cubisch*).

Denn multipliciren wir, um die Bedingungen (2) zum Beweise einführen zu können, $(k\alpha)^3k_x$ mit $(rs)\gtrless 0$, so erhalten wir unter Anwendung des Identitätssatzes:

\noindent\rule{4cm}{0.4pt}

*) Dass die Deutung dieser beiden Gleichungen richtig ist, kann man sofort erkennen, wenn man eine Form $f=a_x^n$ mit Hilfe zweier linearer Formen $\alpha_x$ und $\beta_x$ vermöge der Identität
\[
\alpha_x(\alpha\beta)=(\alpha\beta)\alpha_x-(a\alpha)\beta_x
\]
darzustellen versucht. Erhebt man dieselbe auf die $n^{\text{te}}$ Potenz und entwickelt rechts nach dem binomischen Lehrsatz, so ergiebt sich unmittelbar unter der Voraussetzung $(f,\alpha)^{n-k}=0$, dass $f$ den Factor $\alpha^{k+1}$ besitzt.
%CONT>>


% ===== pdf idx 323, printed 309.0 =====
% !PRINTED 309
\weberpage{309}

\[
\begin{aligned}
(k\alpha)^3k_x(rs) &= (k\alpha)^3\{(ks)r_x-(kr)s_x\} = (k,\alpha^3s)^4-(k,\alpha^3r)^4\\
&= ((k,\alpha s)^2,\alpha^2)^2-((k,\alpha r)^2,\alpha^2)^2=((k,m)^2,\alpha^2)^2-((k,l)^2,\alpha^2)^2\\
&= (n\alpha)^2-(m\alpha)^2=0.
\end{aligned}
\]
Da aber $(rs)\gtrless 0$ ist, so muss $(k\alpha)^3k_x=0$ sein. Es enthält demnach in der That $k$ den Factor $\alpha_x$ zweifach, und folglich nach der Theorie der biquadratischen Formen auch die Hesse'sche Form $\varDelta$ von $k$. Weil sich nun aber alle zweiten Ueberschiebungen von $f$ über $l$, $m$, $n$ als lineare Functionen von $k$, $\varDelta$, $l^2$, $lm$ darstellen lassen, so muss auch jede dieser Formen $(f,l)^2$, $(f,m)^2$, $(f,n)^2$ den linearen Factor $\alpha$ quadratisch enthalten, d.\,h. es muss sein:
\weq{5}{(f,l)^2=a_x^4(a\alpha)(ar)=\alpha_x^2\cdot\varphi}
\weq{6}{(f,m)^2=a_x^4(a\alpha)(as)=\alpha_x^2\cdot\psi,}
wo $\varphi$ und $\psi$ quadratische Formen sind, die für $x=\alpha$ nicht verschwinden. Ueberschiebt man diese beiden Formen dreimal über $\alpha^3$, so müssen diese Ueberschiebungen wegen $(\alpha,\alpha)=0$ verschwinden, d.\,h. man hat:
\[
\begin{aligned}
(\alpha_x^2\cdot\varphi,\alpha_x^3)^3&=a_x(a\alpha)^4(ar)=0\\
(\alpha_x^2\cdot\psi,\alpha_x^3)^3&=a_x(a\alpha)^4(as)=0.
\end{aligned}
\]

Multiplicirt man die erste dieser Gleichungen mit $-s_x$, die zweite mit $r_x$ und addirt, so kommt
\[
a_x(a\alpha)^4\big((as)r_x-(ar)s_x\big)=a_x^2(a\alpha)^4(rs)=0,
\]
oder:
\[
a_x^2(a\alpha)^4=0,
\]
d.\,h. $f$ besitzt den Factor $\alpha_x$ cubisch. Man hat also den Satz:
\glqq Haben die drei quadratischen Formen $l$, $m$, $n$ einen gemeinsamen Factor $\alpha_x$, so besitzt $f$ denselben dreifach.\grqq{}

285. \textit{Umkehrung.} Diese Bedingungen sind nöthig und hinreichend.
Denn besitzt $f$ den Factor $\alpha_x$ dreifach, so ist zunächst
\[
a_x^2(a\alpha)^4=0.
\]
Multipliciren wir anderntheils $(k\alpha)^3k_x$ mit $\alpha_x^4$, so erhalten wir:
\[
(k\alpha)^3k_x\cdot\alpha_x^4=\frac{(ab)^4}{2}\{(a\alpha)^2(b\alpha)b_x+(b\alpha)^2(a\alpha)a_x\}\cdot\alpha_x^4,
\]
oder, weil beide Terme rechts gleich sind:
\[
\begin{aligned}
(k\alpha)^3k_x\cdot\alpha_x^4&=(ab)^4(a\alpha)^2(b\alpha)b_x\cdot\alpha_x^4\\
&=(a\alpha)^2(b\alpha)b_x\cdot\{b_x(a\alpha)-a_x(b\alpha)\}^4.
\end{aligned}
\]
Entwickelt man rechts nach dem binomischen Lehrsatz und führt die Multiplication aus, so besitzt jedes Glied den Factor
\[
(a\alpha)^4a_x^2=(b\alpha)^4b_x^2=0.
\]
%CONT>>


% ===== pdf idx 324, printed 310.0 =====
% !PRINTED 310
\weberpage{310}

%<<CONT
Es verschwindet demnach auch die linke Seite; folglich besitzt $k$ und mit ihm $\varDelta$ den Factor $\alpha_x$ quadratisch. Daraus ergiebt sich endlich unmittelbar, dass
\[
(l\alpha)^2=(m\alpha)^2=(n\alpha)^2=0.
\]
Denn es ist:
\[
(l,\alpha^2)^2=(ak)^4(a\alpha)^2,
\]
also
\[
\alpha^4(l\alpha)^2=(a\alpha)^2\{k_x(a\alpha)-a_x(k\alpha)\}^4=0,
\]
weil
\[
(k\alpha)^3k_x=0,\quad (a\alpha)^4a_x^2=0.
\]
Ferner ist:
\[
(m,\alpha^2)^2=(kl)^2(k\alpha)^2,
\]
also
\[
\alpha^2(m\alpha)^2=(k\alpha)^2\{k_x(l\alpha)-l_x(k\alpha)\}^2=0,
\]
weil
\[
(k\alpha)^3k_x=0,\quad (l\alpha)^2=0.
\]
Endlich ist:
\[
(n\alpha)^2=(km)^2(k\alpha)^2,
\]
also
\[
\alpha^2(n\alpha)^2=(k\alpha)^2\{k_x(m\alpha)-m_x(k\alpha)\}^2=0,
\]
weil
\[
(k\alpha)^3k_x=0,\quad (m\alpha)^2=0.
\]
Die drei quadratischen Formen $l$, $m$, $n$ besitzen also in der That einen gemeinsamen Factor $\alpha$ und es ist demnach $R=0$, $A_{\nu\nu}=0$.

286. \textit{Zweiter Fall; zweite Unterabtheilung.} Die Resultante von $l$ und $m$ verschwindet aber auch, wenn $l$ und $m$ zwei lineare Factoren gemeinschaftlich haben, d.\,h. wenn
\weq{1}{m=\varrho\cdot l,}
wo $\varrho$ eine Proportionalitätsconstante. In diesem Falle ist aber auch $n$ nur bis auf eine Constante von $l$ und $m$ verschieden. Denn
\[
n=(k,m)^2=\varrho(k,l)^2,
\]
und da $(k,l)^2=m$, so hat man:
\[
n=\varrho\cdot m=\varrho^2\cdot l.
\]
Unter diesen Voraussetzungen aber gilt zunächst der Satz:
\glqq Sind die drei quadratischen Covarianten $l$, $m$, $n$ einander proportional und von null verschieden, so kann $A_u$ nicht verschwinden.\grqq{}
Wir beweisen diesen Satz indirect. Wäre nämlich $A_u=0$, also $l$ und damit auch $n$ und $m$ ein volles Quadrat, so hätte man zunächst
\weq{2}{A_u=2C+\frac{AB}{3}=0\ (\text{vgl.\ Nr.\ 267}).}
%CONT>>


% ===== pdf idx 325, printed 311.0 =====
% !PRINTED 311
\weberpage{311}

%<<CONT
Anderntheils war (vgl.\ Nr.\ 266):
\weq{3}{
\left.\begin{aligned}
(f,n)^2&=-\frac{C}{3}k+\frac{B}{3}\varDelta+lm\\
(f,l)^2&=\phantom{-}\frac{A}{3}k+2\varDelta
\end{aligned}\right\}
}
Demnach ist
\[
\begin{aligned}
\Big(f,n-\frac{B}{6}l\Big)^2&=-\frac{C}{3}k+\frac{B}{3}\varDelta+lm-\frac{B}{6}\Big(\frac{A}{3}k+2\varDelta\Big)\\
&=lm-\frac{1}{6}\Big(2C+\frac{AB}{3}\Big)k,
\end{aligned}
\]
oder wegen (2)
\weq{4}{\Big(f,n-\frac{B}{6}l\Big)^2=lm=\varrho l^2.}
Weil aber $n=\varrho^2l$, so ist auch
\weq{5}{\Big(f,n-\frac{B}{6}l\Big)^2=\Big(\varrho^2-\frac{B}{6}\Big)(f,l)^2}
und also durch Comparation
\[
(f,l)^2=\frac{6\varrho}{6\varrho^2-B}l^2.
\]
Nun soll aber $l$ das Quadrat einer linearen Form $\alpha_x$ sein; demnach können wir die letzte Gleichung auch schreiben:
\[
(f,\alpha^2)^2=(a\alpha)^2a_x^4=\frac{6\varrho}{6\varrho^2-B}\alpha_x^4=\text{Const.}\times\alpha_x^4.
\]
Durch einmalige Ueberschiebung dieser Gleichung über $\alpha_x$ folgt:
\[
(a\alpha)^3a_x^3=0,
\]
d.\,h. $f$ besitzt den Factor $\alpha_x$ vierfach.
Dann ist aber auch $k$ das Biquadrat von $\alpha$. Denn in der Gleichung
\[
\alpha^4(k,\alpha)=(ab)^4\alpha^4\cdot a_xb_x^2(a\alpha)=(a\alpha)a_xb_x^2\{(a\alpha)b_x-(b\alpha)a_x\}^4
\]
verschwindet wegen $(a\alpha)^3a_x^3=(b\alpha)^3b_x^3=0$ die ganze rechte Seite und es ist sonach
\[
(k,\alpha)=0,\quad\text{also}\quad k=\text{Const.}\times\alpha^4.
\]
Bilden wir daher die Covariante $l$, also die vierte Ueberschiebung von $f$ über $k$, so ist diese vierte Ueberschiebung null, weil sowohl $f$ als $k$ den linearen Factor $\alpha_x$ vierfach besitzen. Es wäre demnach unter der Annahme $A_u=0$
\[
(f,k)^4=l=0,
\]
das ist aber gegen die Voraussetzung, und sonach muss $A_u\gtrless 0$ sein.

287. \textit{Schärfere Formulirung unserer Voraussetzungen.} Wir können nun unsere Bedingungen schärfer abgrenzen. Sind die drei quadratischen Formen einander proportional, so muss $l$, wegen $A_u\gtrless 0$, zwei von einander verschiedene lineare Factoren besitzen, d.\,h. es muss sein:
%CONT>>


% ===== pdf idx 326, printed 312.0 =====
% !PRINTED 312
\weberpage{312}

%<<CONT
\[
m=\varrho\cdot l=\varrho\cdot\alpha\cdot\beta\quad\text{und}\quad n=\varrho^2\cdot\alpha\cdot\beta,
\]
wo $\alpha$ und $\beta$ die beiden linearen Factoren sind. Es ist demnach
\weq{1}{
\left.\begin{aligned}
(l\alpha)^2&=0,\quad (m\alpha)^2=0,\quad (n\alpha)^2=0\\
(l\beta)^2&=0,\quad (m\beta)^2=0,\quad (n\beta)^2=0.
\end{aligned}\right\}
}
Da aber:
\[
m=\varrho\cdot\alpha\cdot\beta=(k,l)^2=\varrho(k,\alpha\beta)^2
\]
ist, so erhalten wir hieraus:
\weq{$2_{\mathrm{a}}$}{
\left.\begin{aligned}
((k,\alpha\beta)^2,\alpha^2)^2&=(k,\alpha^3\beta)^4=0\\
(k,\alpha\beta^3)^4&=0
\end{aligned}\right\}
}
Ebenso folgen aus diesen Gleichungen (1) für die Form $\varDelta$ die beiden Relationen:
\weq{$2_{\mathrm{b}}$}{
\left.\begin{aligned}
(\varDelta,\alpha^3\beta)&=0\\
(\varDelta,\alpha\beta^3)&=0,
\end{aligned}\right\}
}
wie man erkennt, wenn man diese Ueberschiebungen mit $(\alpha\beta)^2$ multiplicirt, und alsdann $(\alpha\beta)^2(kk_1)^2$ nach dem Identitätssatze entwickelt.
Nun sind die zweiten Ueberschiebungen von $f$ über $l$ und $m$ lineare Functionen von $k$, $\varDelta$ und $l^2$. Wenn wir daher $(f,l)^2$ und $(f,m)^2$ noch viermal über $\alpha^3\cdot\beta$ schieben, so verschwinden die so entstehenden Invarianten, d.\,h. es ist:
\weq{3}{
\left.\begin{aligned}
(f,\alpha^4\beta^2)^6&=0\\
(f,\alpha^2\beta^4)^6&=0
\end{aligned}\right\}
}
In die Relationen $(2_{\mathrm{a}})$ können wir noch die Symbole von $f$ einführen. Es ist:
\[
(k,\alpha^3\beta)^4=(ab)^4(a_x^2b_x^2,\alpha_x^3\beta_x)^4=(ab)^4(a\alpha)^2(b\alpha)(b\beta)=0.
\]
Multipliciren wir diesen Ausdruck mit $(\alpha\beta)^4$, und wenden auf $(ab)^4(\alpha\beta)^4$ den Identitätssatz an, so kommt:
\[
0=(ab)^4(\alpha\beta)^4(a\alpha)^2(b\alpha)(b\beta)\doteq(a\alpha)^2(b\alpha)(b\beta)\{(a\alpha)(b\beta)-(b\alpha)(b\beta)\}^4.
% checked against the GDZ scan (27 Oct 2026): the print has (b\alpha)(b\beta) in the second term (the identity would give (b\alpha)(a\beta)); reproduced as printed
\]
Entwickelt man die rechte Seite nach dem binomischen Lehrsatz, so findet man, dass alle Glieder bis auf das erste wegen der Relationen (3) für sich allein verschwinden. Es muss sonach das erste Glied selbst null sein, d.\,h. wir haben statt der Relationen (1) die beiden folgenden:
\weq{4}{
\left.\begin{aligned}
(a\alpha)^6\cdot(a\alpha)(a\beta)^5&=0\\
(a\beta)^6\cdot(a\beta)(a\alpha)^5&=0,
\end{aligned}\right\}
}
wobei die zweite aus der ersten durch Vertauschung von $\alpha$ mit $\beta$ gewonnen wurde. Hiezu treten die beiden Relationen (3), welche wir schreiben können:
\weq{5}{
\left.\begin{aligned}
(a\alpha)^4(a\beta)^2&=0\\
(a\beta)^4(a\alpha)^2&=0
\end{aligned}\right\}
}


% ===== pdf idx 327, printed 313.0 =====
% !PRINTED 313
\weberpage{313}

Die vier Gleichungen (4) und (5) müssen also zusammen bestehen, wenn $m = \varrho \cdot l = \varrho \cdot \alpha \cdot \beta$ und $n = \varrho^2 \cdot \alpha \cdot \beta$ ist.

Nun kann das erste Gleichungspaar (4) aus drei verschiedenen Ursachen verschwinden. Es kann sein

\weq{1}{(a\alpha)^6=0,\quad (a\beta)(a\alpha)^5=0\ *)}

\weq{2}{(a\alpha)^6=0,\quad (a\beta)^6=0}

\weq{3}{(a\alpha)(a\beta)^5=0,\quad (a\beta)(a\alpha)^5=0.}

Demnach spaltet sich der Fall

\[ n = \varrho\, m = \varrho^2\, l \]

in drei Gruppen, deren jede wir nun einzeln untersuchen wollen.

288. \textit{Erste Gruppe von Bedingungen.} Es sei:

\weq{1}{(a\alpha)^6=0,}

\weq{2}{(a\alpha)^5(a\beta)=0,}

und hiezu

\weq{3}{(a\alpha)^4(a\beta)^2=0,}

\weq{4}{(a\beta)^4(a\alpha)^2=0.}

Sondert man von den ersten drei Bedingungen einen Moment den Factor $(a\alpha)^4$ ab, so sind die übrigbleibenden Factoren

\[ (a\alpha)^2,\quad (a\alpha)(a\beta),\quad (a\beta)^2 \]

die Coefficienten von $\alpha_x$ und $\beta_x$ in der Entwicklung

\[ (\alpha\beta)^2 a_x^2 = \{(a\beta)\alpha_x - (a\alpha)\beta_x\}^2. \]

Multiplicirt man demnach diese Relation wieder mit $(a\alpha)^4$, so kommt:

\[ (a\alpha)^4 a_x^2 (\alpha\beta)^2 = (a\alpha)^4(a\beta)^2\alpha_x^2 - 2(a\alpha)^5(a\alpha) a_x\beta_x + (a\alpha)^6\beta_x^2 = 0. \]

Es ist daher:

\[ (a\alpha)^4 a_x^2 = 0,\ \text{d. h.} \]

\glqq der Factor $\alpha_x$ ist ein cubischer Factor von $f$\grqq{}.

Dieser Fall führt also zu dem nämlichen Resultaten wie der bereits vorhin behandelte: % [sic]

\[ l = \alpha r,\quad m = \alpha s,\quad n = \alpha t. \]

289. \textit{Zweite Gruppe von Bedingungen.} Es möge ferner stattfinden:

\weq{II}{\begin{aligned} (a\alpha)^6&=0, & (a\beta)^6&=0\\ (a\alpha)^4(a\beta)^2&=0, & (a\beta)^4(a\alpha)^2&=0 \end{aligned}}

Versucht man die Form $f$ mit Hilfe der linearen Factoren $\alpha_x$ und $\beta_x$ darzustellen, so erhält man aus der binomischen Entwicklung

\noindent\rule{4cm}{0.4pt}

\noindent *) Der Fall $(a\beta)^6=0,\ (a\alpha)(a\beta)^5=0$ ist nicht verschieden von (1); denn er entsteht durch Vertauschung von $\alpha$ mit $\beta$.
%CONT>>


% ===== pdf idx 328, printed 314.0 =====
% !PRINTED 314
\weberpage{314}

%<<CONT
\[ \{(\alpha\beta) a_x\}^6 = \{(a\beta)\alpha_x - (a\alpha)\beta_x\}^6 \]

unter Anwendung der Relationen (II) das Resultat:

\[ (\alpha\beta)^6 a_x^6 = -\,6(a\beta)(a\alpha)\alpha_x\beta_x \left\{ (a\beta)^4\alpha_x^4 + \frac{10}{3}(a\beta)^2(a\alpha)^2\alpha_x^2\beta_x^2 + (a\alpha)^4\beta_x^4 \right\}, \]

d. h. die Form sechsten Grades zerfällt in das Product einer biquadratischen Form $\Psi$, multiplicirt in die Covariante $l=\alpha\beta$. Die Gleichung sechsten Grades kann also in diesem Falle durch Division mit der bekannten Covariante $l$ auf eine biquadratische reducirt werden. Der biquadratische Factor steht aber noch in engerem Zusammenhange mit $l$. Untersucht man nämlich wie bisher auch die Ueberschiebungen von $k$ über Potenzen von $\alpha$ und $\beta$, so zeigt sich, dass die vierte Ueberschiebung von $k$ über $\alpha^3\beta$ und über $\beta^3\alpha$ verschwindet. Denn man erhält, wenn man um auf die Bedingungen zu gelangen,

\[ (k\alpha)^3(k\beta) = 2(ab)^4(a\alpha)^2(b\alpha)(b\beta) \]

mit $(\alpha\beta)^4$ multiplicirt:

\begin{align*}
(ab)^4(\alpha\beta)^4(a\alpha)^2(b\alpha)(b\beta) &= (a\alpha)^2(b\alpha)(b\beta)\{(b\beta)(a\alpha) - (a\beta)(b\alpha)\}^4\\
&= (a\alpha)^6(b\alpha)(b\beta)^5 - 4(b\beta)^4(b\alpha)^2(a\alpha)^5(a\beta) + 6(a\alpha)^4(a\beta)^2(b\alpha)^3(b\beta)^3\\
&\quad - 4(b\alpha)^4(b\beta)^2(a\alpha)^3(a\beta)^3 + (a\beta)^4(a\alpha)^2(b\alpha)^5(b\beta) = 0.
\end{align*}

Es ist also in der That

\[ (k,\ \alpha^3\beta)^4 = 0 \quad \text{und ebenso} \quad (k,\ \beta^3\alpha)^4=0, \]

und deshalb auch [vgl. Nr. 287 $(2_a)$ und $(2_b)$]

\[ (\varDelta,\ \alpha^3\beta)^4 = 0 \quad \text{und} \quad (\varDelta,\ \beta^3\alpha)^4=0. \]

Will man also $k$ oder $\varDelta$ durch $\alpha$ und $\beta$ darstellen, so verschwinden aus diesem Grunde die Coefficienten der ungeraden Potenzen von $\alpha_x$ und $\beta_x$, d. h. $k$ sowohl als $\varDelta$ werden dadurch in die Normalform transformirt. Demnach sind $\alpha$ und $\beta$ gemäss der Theorie der biquadratischen Formen (vgl. Nr. 179) nichts anderes, als die linearen Factoren einer der drei irrationalen quadratischen Covarianten, in welche die Form $t$ von $k$ sich spalten lässt, d. h.

\glqq\textit{Die Form $f$ ist das Product einer biquadratischen Form in einer ihrer irrationalen quadratischen Covarianten.}\grqq{}

290. \textit{Dritte Gruppe von Bedingungen.} Ist endlich:

\[ (a\alpha)\,(a\beta)^5=0,\quad (a\beta)(a\alpha)^5=0, \]

\[ (a\alpha)^4(a\beta)^2=0,\quad (a\beta)^4(a\alpha)^2=0, \]
%CONT>>


% ===== pdf idx 329, printed 315.0 =====
% !PRINTED 315
\weberpage{315}

%<<CONT
so reducirt sich

\[ \{(\alpha\beta) a_x\}^6 = \{(a\beta)\alpha_x - (a\alpha)\beta_x\}^6 \]

auf die drei Glieder:

\[ (\alpha\beta)^6 a_x^6 = (\alpha\beta)^6\alpha_x^6 - 20(\alpha\beta)^3(a\alpha)^3\alpha_x^3\beta_x^3 + (a\alpha)^6\beta_x^6, \]

d. h. $f=0$ reducirt sich durch die Substitution $\alpha_x^3=x_1$, $\beta_x^3=x_2$ auf eine quadratische Gleichung. Untersucht man überdies die Ueberschiebungen von $k$ über $\alpha$ und $\beta$, so findet sich, dass nicht nur

\[ (k,\ \alpha^3\beta)^4=0 \quad \text{und} \quad (k,\ \alpha\beta^3)^4=0, \]

sondern auch, dass

\[ (k\alpha)^4=0 \quad \text{und} \quad (k\beta)^4=0, \]

wovon man sich leicht auf dem nun schon öfters eingeschlagenen Wege überzeugen kann. Daraus geht aber hervor, dass $k$ das Quadrat von $l$ ist; denn in der binomischen Entwicklung

\[ k_x^4(\alpha\beta)^4 = \{(k\beta)\alpha_x - (k\alpha)\beta_x\}^4 \]

verschwinden alle Glieder bis auf das Glied $(k\alpha)^2(k\beta)^2\alpha_x^2\beta_x^2$.

291. \textit{Zusammenfassung der Resultate.} Wir haben demnach folgende Sätze: Ist

\[ m = \varrho \cdot l = \varrho \cdot \alpha \cdot \beta, \]

so ist

\begin{enumerate}
\item[1)] entweder $\alpha_x$ ein cubischer Factor von $f$, und $\alpha_x$ ergiebt sich rational nach der Methode für Aufsuchung des grössten gemeinschaftlichen Theilers von $f$ und $l$,
\item[2)] oder $f$ zerfällt in eine biquadratische Form und die Covariante $l$, welche zugleich Factor der Form $t$ von jener Form $k$ ist,
\item[3)] oder endlich $f$ lässt sich durch die Substitution
\end{enumerate}

\[ x_1 = \alpha_x^3,\quad x_2 = \beta_x^3 \]

auf eine quadratische Gleichung reduciren.

292. \textit{Zweiter Fall; dritte Unterabtheilung.} Die Resultante $A_{\nu\nu}$ verschwindet endlich auch, wenn eine der quadratischen Formen $l$ oder $m$ null ist. Verschwindet $l$, so ist $m$ a priori null, da ja $m=(k,l)^2$ ist. Man kann aber auch beweisen, dass, wenn $m=0$, nothwendig auch $l$ verschwinden muss. Nehmen wir nämlich an, $l$ sei $\gtrless 0$, wenn $m=0$, dann ist wegen

\[ 0 = (\varDelta,\ m)^2 = \frac{B}{6}\, m + \frac{C}{3}\, l \quad \text{(vgl. Nr. 265, 2)}, \]

\[ \frac{C}{3}\, l = 0,\quad \text{also}\quad C=0, \]

und demnach, da ja $A_{lm}=(l,m)^2=0=\dfrac{2}{3}(B^2+AC)$ auch $B=0$.


% ===== pdf idx 330, printed 316.0 =====
% !PRINTED 316
\weberpage{316}

Es ist aber

\[ 2C + \frac{AB}{3} = A_{ll} \]

und somit verschwindet auch $A_{ll}$.

Nun ist (vgl. Nr. 266, 5)

\[ 0 = (f,\ m)^2 = \frac{l^2}{2} + \frac{B}{3}\, k + \frac{A}{3}\, \varDelta, \]

also

\[ 0 = \frac{l^2}{2} + \frac{A}{3}\, \varDelta. \]

Schieben wir anderntheils $(f,\ m)^2$ zweimal über $f$, so erhalten wir % checked against the GDZ scan (27 Oct 2026): the print reads "anderntheils"

\begin{align*}
\big((f,\ m)^2,\ f\big)^2 &= \frac{1}{2}(l^2,\ f)^2 + \frac{A}{3}(\varDelta,\ f)^2\\
&= \frac{1}{2}\left\{\Big(2\varDelta + \frac{A}{3}\, k\Big) l - \frac{1}{3} A_{ll} f\right\} + \frac{A}{3}\left\{\frac{lk}{2} + \frac{B}{6} f\right\}\\
&= l\varDelta + \frac{lkA}{3}.
\end{align*}

Dieser Ausdruck ist null, da $(f,\ m)^2$ verschwindet; da aber $l$ nach Annahme von null verschieden sein soll, so ist

\[ \varDelta + \frac{A}{3}\, k = 0. \]

Schieben wir diesen Ausdruck zweimal über $k$, so kommt

\[ (\varDelta,\ k)^2 + \frac{A}{3}\, \varDelta = 0 = \frac{B}{6}\, f + \frac{A}{3}\, \varDelta. \]

Daher ist:

\[ \frac{A}{3}\, \varDelta = 0. \]

Es müsste also auch wegen

\[ \frac{l^2}{2} + \frac{A}{3}\, \varDelta = 0 \]

$l^2=0$ sein, was gegen unsere Annahme ist.

Wenn also eine der drei quadratischen Formen verschwindet, so verschwinden auch die beiden andern.

Nehmen wir nun an

\[ (1)\ l=0,\quad (2)\ k \lessgtr 0. \]

Dann ist, wegen $(f,\ l)^2 = 2\varDelta + \dfrac{A}{3} k = 0$

\[ \varDelta = \varrho \cdot k, \]

d. h. nach der Theorie der biquadratischen Formen: $k$ ist entweder ein volles Quadrat, etwa

\[ k = \alpha^2 \cdot \beta^2 \]
%CONT>>


% ===== pdf idx 331, printed 317.0 =====
% !PRINTED 317
\weberpage{317}

oder auch $k$ ist ein Biquadrat, etwa

\[
k = \alpha^4.
\]

Im ersten Falle, in welchem $(\alpha\beta)\genfrac{}{}{0pt}{}{\leqq}{\geqq}0$, ist also nach Voraussetzung

\[
l = (f, k)^4 = 0.
\]

Weil aber nach Nr. 253 (1) in jedem Falle $(f, k)^3$ gleich null ist, so hat man auch:

\[
(ak)^3 a_x^3 k_y = 0,
\]

da die Elementarcovarianten rechts verschwinden. Ersetzen wir links $k$ durch $\alpha^2\beta^2$, so erhalten wir

\[
a_x^3 (a\alpha)(a\beta)\{\alpha_y(a\beta) + \beta_y(a\alpha)\} = 0.
\]

Nun sind $\alpha$ und $\beta$ von einander verschieden und da diese Relation für jeden Werth von $y$ gilt, so müssen die Coefficienten von $\alpha_y$ und $\beta_y$ einzeln verschwinden, d. h. es muss sein:

\[
\begin{aligned}
a_x^3(a\alpha)(a\beta)^2 &= 0\\
a_x^3(a\alpha)^2(a\beta) &= 0.
\end{aligned}
\]

Multipliciren wir die erste Gleichung mit $\alpha_x$, die zweite mit $-\beta_x$ und addiren, so kommt:

\[
a_x^3(a\alpha)(a\beta)\{(a\beta)\alpha_x - (\alpha a)\beta_x\} = a_x^4(a\alpha)(a\beta)(\alpha\beta) = 0,
\]

also:

\[
a_x^4(a\alpha)(a\beta) = 0.
\]

Versucht man daher $a_x^6$ vermittelst der Identität

\[
a_x^6(\alpha\beta)^6 = \{(\alpha\beta)\alpha_x - (a\alpha)\beta_x\}^6
\]

darzustellen, so reducirt sich die binomische Entwicklung in diesem Falle auf

\[
a_x^6(a\beta)^6 = (a\beta)^6\alpha_x^6 - (a\alpha)^6\beta_x^6,
\]

d. h. \glqq Die Form $f$ reducirt sich auf eine binomische Gleichung, wenn $R=0$, $l=0$, $k=\alpha^2\beta^2$.\grqq{}

Ist aber neben $l=0$ die Form $k$ ein Biquadrat, also

\weq{1}{k = \alpha^4,}

dann ist immer noch:

\[
(f, k)^3 = 0,
\]

also wegen (1)

\[
a_x^3(a\alpha)^3 = 0,
\]

d. h. \glqq Die Form $f$ besitzt den linearen Factor $\alpha$ vierfach.\grqq{}

Umgekehrt: Besitzt $f$ einen vierfachen linearen Factor $\alpha$, so ist $k$ ein Biquadrat und $l$ identisch null. Denn es ist dann:
%CONT>>


% ===== pdf idx 332, printed 318.0 =====
% !PRINTED 318
\weberpage{318}
%<<CONT

\[
\begin{aligned}
\alpha^4(k\alpha) &= (ab)^4\alpha^4\cdot(a\alpha)\alpha_x\beta_x^2\\
&= (a\alpha)a_xb_x^2\{(a\alpha)b_x - (b\alpha)a_x\}^4\\
&= 0, \quad \text{wegen } a_x^3(a\alpha)^3 = 0.
\end{aligned}
\]

Also: $(k, \alpha) = 0$, d. h. $k = \alpha_x^4$.

Deshalb ist aber auch $(f, k)^4 = 0$, d. h. $l = 0$, wie behauptet wurde.

\textit{Anmerkung.} Ist $k$ identisch null, dann ist $f$ entweder eine Octaederform und kann nach der in Nr. 195 gegebenen Methode aufgelöst werden, oder $f$ besitzt einen linearen Factor $\alpha_x$ mindestens fünffach [vgl. § 19, über die Formen, für welche $(f,f)^4=0$].

293. \textit{Dritter Fall}: $A=0$, $C=0$. Wir wollen endlich noch zwei Fälle erwähnen, in welchen die Invariantenrelationen zwar nicht auf algebraisch, so doch mit Hilfe elliptischer Functionen lösbare Gleichungen herführen. (Vgl. Clebsch § 114 der \glqq Binären Formen\grqq{}.) Bildet man nämlich das Formensystem der Modulargleichung (Jacobi, \glqq Fundamenta\grqq{}, S. 27) mit dem einen Parameter $u$:

\weq{1}{v^6 + 4u^5v^5 + 5u^2v^4 - 5u^4v^2 - 4uv - u^6 = 0}

für die Transformation fünfter Ordnung*), so zeigt sich, dass die Invarianten $A$ und $C$ identisch verschwinden. Umgekehrt hat dann Clebsch (a. a. O.) gezeigt, dass $f = a_x^6$ sich durch eine lineare Transformation in die Form

\weq{2}{\varphi = \xi^6 + 5\xi^4\eta^2 + 15\xi^2\eta^4 - 4Z\xi\eta^5 - 5\eta^6}

mit dem einen Parameter $Z$, welche aus (1) durch die Substitution

\[
v = u\,\frac{\dfrac{\xi}{\eta} + u^4}{1 - u^4\dfrac{\xi}{\eta}}
\]

hervorgeht, stets überführen lässt, sobald ihre Invarianten $A$ und $C$ verschwinden. In diesem Falle besitzen die beiden quadratischen Covarianten $l$ und $\mu$ einen gemeinsamen Factor und

\[
\frac{\eta}{\xi} = \frac{l\sqrt{-BD}}{2\mu}
\]

ist die erwähnte lineare Transformation der Form $f$ in $\varphi$.

\footnote{*) Zur Erklärung möge hier beigefügt werden: Ist $k$ der Modul des elliptischen Integrals $\displaystyle\int\frac{dx}{\sqrt{(1-x^2)(1-k^2x^2)}}$, $\lambda$ der Modul des neuen Integrals, das aus ihm durch eine Transformation fünfter Ordnung hervorgeht, dann besteht zwischen $\sqrt[4]{k}=u$ und $\sqrt[4]{\lambda}=v$ die Gleichung (1).}


% ===== pdf idx 333, printed 319.0 =====
% !PRINTED 319
\weberpage{319}

294. \textit{Vierter Fall}: $B=\dfrac{7}{50}A^2$, $C=-\dfrac{9}{500}A^3$. Eine zweite Gleichungsgattung mit einem einzigen Parameter repräsentirt die Multiplicatorgleichung von Brioschi:

\weq{3}{z^6 - 4z^5 + 256k^2k'^2(z+1) = 0, \quad (k'^2 = 1 - k^2),}

auf welche die Frage nach dem Multiplicator $M$ führt, vermöge welchem die Beziehung besteht:

\[
\frac{dy}{\sqrt{(1-y^2)(1-\lambda^2y^2)}} = \frac{1}{M}\cdot\frac{dx}{\sqrt{(1-x^2)(1-\varkappa^2x^2)}}.
\]

Dieser Multiplicator $M$ hängt mit $v$ in (1) durch die Gleichung

\[
M = \frac{v(1-uv^3)}{v-u^5}
\]

und mit $z$ in (3) durch

\[
M = \frac{1}{z+1}
\]

zusammen. Setzt man in (3) $z=\lambda\cdot\dfrac{\xi}{\eta}$, so kann man sie nach Clebsch in der Form schreiben:

\weq{4}{a_0\xi^6 + 6a_1\xi^5\eta + 6a_5\xi\eta^5 + a_6\eta^6 = 0,}

wobei

\[
a_0 = \varrho\cdot\lambda^6, \quad a_1 = -\frac{2}{3}\varrho\lambda^5,
\]

\[
a_5 = \frac{128}{3}\varrho\cdot\lambda k^2k'^2, \quad a_6 = 256\varrho\cdot k^2k'^2
\]

und

\[
a_0a_6 + 9a_1a_5 = 0.
\]

Für diese Gleichung (4) findet man nun durch directe Rechnung die beiden sie charakterisirenden Invarianteneigenschaften:

\[
B = \frac{7}{50}A^2, \quad C = -\frac{9}{500}A^3,
\]

und umgekehrt sind dieselben genügend, um die allgemeine Form $f$ durch die aus

\[
-8(a_1^3a_6^2 + a_0^2a_5^3)\xi\cdot\eta = m - \frac{A}{10}l
\]

sich ergebende Transformation in die Form (4) überzuführen.

\textit{Anmerkung 1.} Die Modulargleichung (1) wurde von Hermite zur Auflösung der Gleichung fünften Grades benutzt. Er fand, dass diese Gleichung sechsten Grades eine Resolvente fünften Grades besitzt, für welche $C=0$ ist, die sich also in die Form bringen lässt:

\[
t^5 - t - \frac{1}{5\sqrt[4]{z}} = 0.
\]


% ===== pdf idx 334, printed 320.0 =====
% !PRINTED 320
\weberpage{320}

Brioschi dagegen ging von der Multiplicatorgleichung (3) aus und kam von dieser auf dieselbe Resolvente wie Hermite. (Vgl. die diesbezüglichen interessanten historischen Entwicklungen in: Klein, \glqq Vorlesungen über das Ikosaeder\grqq{}, insbesondere § 3, S. 144.)

\textit{Anmerkung 2.} Ich erwähne hier noch der Vollständigkeit halber die beiden Brill'schen Normalformen der Gleichung sechsten Grades

\[
\begin{aligned}
&x^6 + 2px^5 + 3qx^4 + 4rx^3 + 3x^2 + 2px + q = 0\\
&x^6 + ax^4 + bx^3 + cx^2 + 1 = 0,
\end{aligned}
\]

in welche sich jede Gleichung sechsten Grades nach Adjunction einer Wurzel einer Gleichung fünften Grades (vgl. Nr. 273--276) linear transformiren lässt. Die Discriminante dieser Gleichungen zerfällt in zwei rationale Factoren (vgl. Math. Annalen, Bd. XX, S. 330).

\begin{center}
\S~31. \gesperrt{Simultanes System einer quadratischen und cubischen Form.}
\end{center}

295. \textit{Allgemeine Betrachtungen.} Die simultanen Systeme mehrerer binärer Formen spielen in vielen Problemen der Mathematik eine hervorragende Rolle. Ich erinnere nur an die Lösung der Gleichung neunten Grades, welcher die neun Wendepunkte einer Curve dritter Ordnung genügen, an die Dreitheilung elliptischer und hyperelliptischer Functionen, an die Reduction elliptischer Integrale erster Gattung auf die Normalform, an die Theorie rationaler Curven, an die Untersuchungen über apolare und involutorische Systeme etc. Es erscheint daher angezeigt, an Beispielen zu zeigen, in welcher Weise derartige Systeme aufgestellt werden. Die einfachsten und sofort allgemein zu behandelnden derartigen Systeme sind simultane Systeme von $n$ linearen Formen, die aus den $n$ Formen selbst und den $\dfrac{n(n-1)}{2}$ Determinanten $(ab)$ bestehen; ferner das simultane System von $n$ quadratischen Formen, das wir bereits früher aufgestellt haben. Es bietet auch wenig Schwierigkeiten, die Formen des vollständigen Systemes zu ermitteln, das entsteht, wenn ein bereits bekanntes volles System mit dem vollen System von $n$ linearen oder $n$ quadratischen Formen überschoben werden soll. Dagegen werden die Untersuchungen bereits complicirter, sobald es sich auch nur um das simultane System von zwei Formen handelt, von denen keine quadratisch oder linear ist. Wir werden demnach im Folgenden zwei Beispiele ins Auge fassen; in dem ersten wird eine der Grundformen eine quadratische Form sein, die andere cubisch; im zweiten werden beide Formen den dritten Grad besitzen. Wir wenden uns nun dazu, das simultane System einer quadratischen und cubischen Form zu studiren.


% ===== pdf idx 335, printed 321.0 =====
